% ============================================================
%  example.tex  –  demonstrates usage of kimletter.sty
%  Compile with:  pdflatex example.tex
% ============================================================
\documentclass[11pt, a4paper]{letter}
\usepackage{letter-mpiea-kim}

% -------------------------------------------------------
%  Override any default fields here (all are optional)
% -------------------------------------------------------
% \klset{name}      {Dr. Seung-Goo Kim}
% \klset{title}     {Postdoctoral Researcher}
% \klset{group}     {Research Group Neurocognition of Music and Language}
% \klset{institute} {MAX PLANCK INSTITUTE\\FOR EMPIRICAL AESTHETICS}
% \klset{postal}    {Grüneburgweg 14, DE-60322 Frankfurt am Main}
% \klset{phone}     {+49\,(0)\,69\,8300\,479\,851}
% \klset{email}     {seung-goo.kim@ae.mpg.de}
% \klset{web}       {https://seunggookim.github.io}
% \klset{city}      {Frankfurt am Main}
% \klset{logo}      {mpiea_logo.png}   % uncomment + set path to show logo

\begin{document}

% --- Letterhead + date ---
\letterhead

% --- Recipient ---
\begin{flushleft}
  \small
  The Editorial Office \\
  \textbf{Annals of the New York Academy of Sciences}\\
  7 World Trade Center, 40th Floor \\
  New York, NY 10007-2157, USA
  \vspace{1em}
\end{flushleft}

% --- Subject line ---
\lettsubject{Submission of Original Resarch Article}

% --- Salutation ---
Dear Editor-in-Chief,
\vspace{0.5em}

% --- Body ---
We are pleased to submit our original research article, titled \emph{"Temporal Structure Perception in Tonal and Atonal Music: An Online Behavioural Study"}, for consideration for publication in the Annals of the New York Academy of Sciences.  This work investigates a fundamental question in music cognition: how does tonal structure shape the perception of temporal information in music? 
% We believe the article is a strong fit for the journal given its integrative scope, empirical rigour, and relevance to the neuroscience and psychology of music.

Across five online behavioural experiments (N = 487), we used a `quilting' algorithm to rearrange segments of tonal (Bach) and atonal (Krenek) violin recordings across a wide range of temporal scales (60--3840 ms), and asked participants to rate perceived naturalness. Naturalness increased with segment length for both corpora, following a piecewise linear function with comparable slopes and elbow points, but a consistently higher plateau for tonal music, indicating that tonal and temporal structures contribute independently and additively to perception. A cross-domain comparison with speech revealed a reversed familiarity effect, suggesting domain-specificity in musical temporal processing. Musicians showed heightened sensitivity specifically in atonal music, while in non-musicians, self-reported Perceptual Abilities (Gold-MSI) robustly predicted the tonal elbow point and the tonal-atonal plateau difference, replicated via cross-validation.

These findings speak directly to the journal's interdisciplinary scope at the intersection of music psychology and auditory neuroscience. We confirm the manuscript is original, not under consideration elsewhere, and that all authors have approved the submission. The study was conducted in accordance with a protocol approved by the Duke University Institutional Review Board (IRB protocol number: D0206). Informed consent was obtained from all participants. The authors declare no conflicts of interest.

% To address this question, we conducted five preregistered online behavioural experiments with 487 participants. Drawing on a dedicated ‘quilting’ algorithm, we systematically rearranged segments of tonal (Bach) and atonal (Krenek) violin recordings while preserving temporal structures spanning 60 to 3840 ms. Participants rated the perceived naturalness of each excerpt, allowing us to quantify how segment duration and tonal context jointly determine perceptual coherence across a broad and ecologically representative sample.

% Naturalness ratings increased with intact segment length for both corpora, following a piecewise linear function with comparable slopes and elbow points—but a consistently higher plateau for tonal music. This pattern suggests that tonal and temporal structures contribute independently and additively to naturalness perception rather than interacting. A cross-domain comparison with speech quilts revealed a reversed familiarity effect, pointing to domain-specificity in musical temporal processing. Furthermore, musicians showed heightened sensitivity specifically in atonal music, while in non-musicians, self-reported Perceptual Abilities (Gold-MSI) were robustly associated with the tonal music elbow point and the tonal–atonal plateau difference—a relationship replicated via cross-validation.

% These findings advance our understanding of how musical expertise and innate musicality independently modulate temporal integration in music, demonstrating that parallel processing of tonal and temporal structure is a core feature of music perception. The results carry implications for theories of auditory scene analysis, predictive coding in music, and the neural bases of rhythm and meter. The study also contributes methodologically by demonstrating the scalability and validity of online quilting paradigms for investigating auditory temporal processing.

% The Annals of the New York Academy of Sciences has long served as a premier venue for interdisciplinary research on music, cognition, and the brain, including landmark volumes on the neurosciences and music. Our manuscript engages directly with this tradition by combining carefully controlled behavioural experimentation with theoretical perspectives drawn from music psychology, auditory neuroscience, and computational modelling. We are confident the article will resonate with the journal’s broad and distinguished readership.

% We confirm that this manuscript is original, has not been published previously, and is not under consideration for publication elsewhere. All authors have approved the final version of the manuscript and agree to its submission. The authors declare no conflicts of interest. 

% We would be happy to suggest potential reviewers or provide any additional information you may require. We look forward to hearing from you at your earliest convenience and thank you sincerely for your time and consideration.


% --- Sign-off (optional argument sets signature gap height) ---
\lettsignoff[3em]

On behalf of all co-authors. %: \\
% Dr. Daniel Müllensiefen, University of Hamburg, Hamburg, Germany\\
% Ms. Ying Yu, Columbia University Vagelos College, New York, U.S.A.\\
% Dr. Tobias Overath, Duke University, North Carolina, U.S.A.

\end{document}
